% ============================================================
%  Section 7 — Conclusion  (~0.5 pages)
% ============================================================
\section{Conclusion}

We introduced a unified benchmarking methodology in which QCanvas
compiles circuits from Qiskit, Cirq, and PennyLane to a common OpenQASM
3.0 intermediate representation for execution on a single QSim
statevector backend, eliminating the simulator confound present in all
prior cross-framework comparisons.
Across 12 algorithms and 6 metric dimensions, Qiskit and Cirq prove
functionally equivalent ($S_{\text{overall}}$ 58.80 vs.\ 58.79), with
selection between them a matter of ecosystem fit rather than performance.
PennyLane exhibits initial super-linear gate scaling for oracle and search
algorithms, with Grover's Algorithm reaching 216 gates at 6 qubits
(a $12\times$ overhead relative to the 18 gates produced by Qiskit and
Cirq), though this transitions to a linear growth regime ($a=0.218$) for
qubits $n \ge 7$.
Most critically, PennyLane's gate decomposition produces JSD$=1.000$
divergence against both other frameworks for six algorithms at 4 or more
qubits, a systematic incompatibility in measurement basis that is
invisible without a common simulation backend.
Future work includes physical hardware validation on NISQ devices,
extension of the suite to fault-tolerant algorithms, and integration of
additional frameworks including Amazon Braket and Stim.
